\documentclass[]{article}

\usepackage{layout}
\usepackage{amsmath,amssymb}
\usepackage{booktabs}
\usepackage{amsthm}
\usepackage{appendix}
\usepackage{graphicx}

\usepackage[backend=biber,
style=authoryear,
sorting=nyt,
uniquename=false
]{biblatex}
\addbibresource{macrotest.bib}
\bibliography{macrotest}
%\theoremstyle{plain}% Theorem-like structures provided by amsthm.sty
\newtheorem{theorem}{Theorem}[section]
\newtheorem{lemma}[theorem]{Lemma}
\newtheorem{corollary}[theorem]{Corollary}
\newtheorem{proposition}[theorem]{Proposition}

%\theoremstyle{definition}
\newtheorem{definition}[theorem]{Definition}
\newtheorem{example}[theorem]{Example}
\newtheorem{result}{Result}

%\theoremstyle{remark}
\newtheorem{remark}{Remark}
\newtheorem{notation}{Notation}
\newtheorem{assumption}{Assumption}

\usepackage{geometry}
\usepackage{setspace}
\usepackage{multirow}
\usepackage{hhline}
%\usepackage{cellspace}

\usepackage{subcaption}
%\setlength\cellspacetoplimit{5pt}            %Using booktabs make this superfluous?
%\setlength\cellspacebottomlimit{5pt}
\geometry{lmargin=1.5in,rmargin=1.5in}

\begin{document}

\title{Shotgun Wedding: The Unpleasant Monetarist Arithmetic and Modern Monetary Theory}
\date{\textsf{\today}}
\author{George A. Waters\footnote{email:  gawater@gmail.com;  website:  about.illinoisstate.edu/gawater}\footnote{There are no financial relationships or financial support to report.}
\\Illinois State University}
\maketitle


\begin{abstract}
Stabilization policy conducted through the interest rate on bank reserves and partial monetary finance of fiscal deficits is a potentially beneficial policy combination. Under the \textit{unpleasant monetarist arithmetic}, a sale of bonds by the monetary authority raises steady state inflation due to the higher interest cost of the net debt. The same intuition suggests that monetary finance of fiscal deficits, the purchase of bonds by the monetary authority, as advocated by Modern Monetary Theorists, would lower inflation. While full monetary finance is not recommended, as it eliminates the bond market, partial monetary finance with monetary policy targeting interest on reserves could be beneficial for countries with existing debt in excess of the needs of financial markets.

This paper presents a shopping time model of money demand with a banking sector where deposits provide liquidity services but require bank reserves. Fiscal and monetary budget constraints are treated separately to allow for the study of monetary finance. The resulting consolidated budget constraint demonstrates that monetary seigniorage must cover the cost of the fiscal deficit, the interest on the debt, and the interest on reserves. Higher costs imply higher steady state inflation, and the model embodies the intuition about monetary finance and the unpleasant monetarist arithmetic. Alternatively, for a given inflation target there is a trade-off between the fiscal deficit and the steady state interest rate on reserves. A calibrated version of the model verifies these intuitions, shows that the underlying dynamic model is determinate and expectationally stable when the interest rate rule for reserves satisfies the Taylor Principle. Bayesian estimation of an extension of the model shows that U.S. policy satisfies the stability conditions for the sample where monetary policy targets the interest rate on reserves.

\textbf{Keywords: \ }Unpleasant Monetarist
Arithmetic, Monetary Finance, Fiscal-Monetary models, Expectational Stability, Modern Monetary Theory

\textbf{JEL classification: }A22, E0, E52, E60
\end{abstract}


\pagebreak

\onehalfspacing
\section{Introduction}
 
The \textit{unpleasant monetarist arithmetic} (\cite{sargent1981unpleasant}) conveys the idea that a sale of bonds by the monetary authority can increase inflation due to the higher interest costs of the net debt. If selling bonds is bad, isn't buying bonds good? Modern Monetary Theorists (MMT), such as \textcite{Kelton2020}, advocate for full monetary finance, equivalent to the monetary authority buying all government debt. 

Monetary finance has a bad reputation, as it is the cause of most hyperinflations, \textcite{sargent1982ends}. However, Japan is notable case where monetary finance has been used responsibly. One objection, particularly relevant before the Great Recession, is that full monetary finance impedes the conduct of open market operations, since it eliminates the bond market. Furthermore, \textcite{AIYAGARI98debt} and others\footnote{\textcite{Woodford90debt}, \textcite{Azzimonti2019debtwithbanks}, and \textcite{brunnermeier24safe} have similar implications for optimal debt.} demonstrate that government debt held by the public alleviates financial market frictions, so full monetary finance cannot be recommended.

However, the use of partial monetary finance while targeting the interest rate on reserves for short run stabilization is a policy combination that merits consideration, particularly for countries with existing debt in excess of that required by financial markets. Monetary finance reduces the interest payments on the net debt, though the interest on reserves constitutes an additional cost. 

The present work combines a shopping time model of money demand with a banking sector that includes reserves, while accounting for the fiscal and monetary budget constraints separately to allow for the study of monetary finance. The resulting steady state consolidated government budget constraint contains much of the intuition. Higher costs to the government, including fiscal deficits, interest on the debt, and interest on reserves, requires higher inflation to generate the necessary seigniorage.

The unpleasant monetarist arithmetic and the benefit of monetary finance are both apparent. The greater the fraction of debt purchased by the monetary authority, the lower the steady state inflation rate and vice versa. From another perspective, for a given inflation target and degree of monetary finance, there is a trade-off between the fiscal deficit and the interest on reserves. The higher the deficit, the lower the steady state interest rate on reserves must be to achieve the steady state inflation target.

For the underlying dynamic model with expectations, active monetary policy is conducted through the interest rate on reserves, and active fiscal policy adjusts taxes in response the real net debt. Active fiscal policy plays a minimal role in achieving determinacy and expectational stability. In fact, overly aggressive fiscal policy can be destabilizing. Determinacy and expectational stability are achieved with an interest rate rule where the interest rate on reserves satisfies the Taylor principle, similar to the results in \textcite{bullardmitra02}. \textcite{eusepi2018fiscal} study a model with active fiscal and monetary policy, though without a financial sector. The monetary policy requirements for long run expectational stability are similar. Bayesian estimation of an extension of the model shows that U.S. policy satisfies the stability conditions for the sample 2009-2025 when monetary policy targets the interest rate on reserves.

For the short term conduct of monetary policy, there is an inverse relationship between policy rates and subsequent inflation. In the steady state analysis, however, higher costs to the government cause inflation and nominal rates to rise together, as with the Neo-Fisherian policy in \textcite{cochrane2023NeoFisher}. The contrast between the short and long run intuitions echoes the interpretations of open market operations in the first paragraph. An open market sale of bonds increases the interest on the net debt, which leads to higher inflation in the steady state analysis. Such a sale is permanent. In the short run, an open market sale to lower the money supply can reduce inflation due to liquidity effects, but it is likely followed by a bond purchase in the future.

There are no countries where the fraction of bonds purchased by the monetary authority is an explicit target, but some countries do use monetary finance. The European Central Bank has made indirect purchases of government bonds to control spikes in yields of the bonds for some countries, as discussed in \textcite{Attavila2015OMT}. The Bank of Japan currently holds 43\% of the debt issued by the fiscal authorities, which helps to explain the limited impact of a gross debt well over 200\% of GDP. Since 2010, Japanese policy has been \textit{de facto} close to the policy proposed here. The net debt has stabilized around 120\% of GDP, largely due to bond purchases by the Bank of Japan, and monetary policy is conducted using the interest rate on reserves. See \textcite{Hansen2023JapanDebt} and \textcite{waters2024potentiallypleasant} for more detailed discussions.

The model builds on those of \textcite{ireland2014macroeconomics} and Chapter 27 of \textcite{sargent2018recursive} who both develop shopping time models of money demand. The former includes a banking sector to study monetary policy conducted through interest on reserves, which uses a model of deposits based on \textcite{christiano1995inside}. \textcite{Belongia2014monies} use a similar model to study the relationships between monetary aggregates. \textcite{sargent2018recursive} focuses on fiscal-monetary interactions including the unpleasant monetarist arithmetic. \textcite{bassetto2013fiscal} develop a model treating fiscal and monetary policy constraints independently to determine conditions for central bank profitability. 

Models of the fiscal theory of the price level (\cite{leith2016understanding}, \cite{cochrane2022ftpl}) show interactions of fiscal and monetary policy and their impact on prices as well. \textcite{bassettosargent2020shotgun} discuss the interpretation of the unpleasant monetarist arithmetic, the fiscal theory of the price level and the quantity theory of money in a single model. 

The paper is organized as follows. Section 2 derives the consolidated budget constraint and discusses the results from the steady state version. Section 3 describes the micro-founded model with a banking sector that determines the demand for money and reserves. Section 4 details a calibrated version of the model with results on determinacy and expectational stability, as well as quantitative results about the policy implications. Section 5 reports Bayesian estimation results, and Section 6 concludes.

\section{Budget Constraints}

To examine monetary financing of fiscal deficits and interest on reserves together,
consider the flow budget constraints for fiscal and monetary policy
separately. The fiscal constraint is%
\begin{equation}
 P_{t}g_{t}=P_{t}\tau _{t}+r_{t}^{-1}B_{t}^{T}-B_{t-1}^{T}+RCB_{t}
 \label{fiscal BC}
\end{equation}
where the price level is $P_{t},$ and real government spending and taxes are $%
g_{t}$ and $\tau _{t}$. One period discount bonds $B_{t}^{T}$ are issued by the fiscal authority and paid off at a gross
yield of $r_{t}$. The fiscal authority also receives a transfer $RCB_{t}$
from the monetary authority.

The monetary authority pays gross interest $r_{t}^{n}$ on reserves $N_{t}$ and has the flow budget constraint%
\begin{equation}
 r_{t}^{-1}B_{t}^{M}-B_{t-1}^{M}+RCB_{t}=M_{t}-M_{t-1}-(r_{t}^{n}-1)N_{t}
 \label{CB BC}
\end{equation}
where bonds held by the central bank is $B_{t}^{M}$, and currency is $M_{t}$. The timing of the $N_{t}$ term reflects the short term payments of interest on reserve deposits and matches the expression for bank profits.

Monetary financing is represented by the fraction $\Omega$ of the bonds issued that are purchased by the monetary authority. The bonds held by the public or net debt is defined as
\[
B_{t}=B_{t}^{T}-B_{t}^{M},
\]
so the following hold.
\begin{eqnarray}
 B_{t}^{M}=\Omega B_{t}^{T} \\
 B_{t}=(1-\Omega) B_{t}^{T} \label{net debt}
\end{eqnarray}

Using the above relation for $B_{t}^{M}$ and substituting for $RCB_{t}$ in the above fiscal and monetary constraints yields the following consolidated budget constraint.%
\begin{equation}
M_{t}=M_{t-1}+P_{t}g_{t}-P_{t}\tau _{t}+(r_{t}^{n}-1)N_{t}-(1-\Omega)(r_{t}^{-1}B_{t}^{T}-B_{t-1}^{T})
\label{cons bc}
\end{equation}

The MMT case of full monetary finance is equivalent to the monetary authority purchasing all the bonds issued so $\Omega=1$ and $B_{t}^{M}=B_{t}^{T}$. Partial monetary finance means $\Omega\in(0,1)$.

The policy of purchasing bonds is a money supply rule. Given all the other values including the fraction purchased $\Omega$, the monetary authority sets $M_{t}$ according to the constraint (\ref{cons bc}). However, active monetary policy can still be conducted via interest on reserves. Let the interest rate rule for $r^{n}_{t}$ depend on gross inflation
\[
\Pi_{t}=\frac{P_{t}}{P_{t-1}}
\]
relative to steady state inflation $\Pi$ as follows.
\begin{equation} \label{ior}
    r^{n}_{t}=r_{t}^{\alpha_{i}}\left(\frac{\Pi_{t}}{\Pi}\right)^{b_{\pi}}
\end{equation}
The interest rate is set below the market rate on bonds so the parameter $\alpha_{i}$ is such that $0\leq\alpha_{i}<1$. The policy rate is set to respond directly to inflation so the parameter $b_{\pi}$ is non-negative, and the Taylor principle is satisfied when $b_{\pi}>1$. 

\subsection{Real steady state consolidated budget constraint}

The consolidated budget constraint can be written in real terms where real currency and real bonds are
\[m_{t}=\frac{M_{t}}{P_{t}},\] and \[b_{t}=\frac{B_{t}}{P_{t}}.\]

The constraint becomes
\begin{equation}
g_{t}-\tau _{t}+(r_{t}^{n}-1)n_{t}=(1-\Omega)(r_{t}^{-1}b_{t}^{T}-\Pi^{-1}_{t}b_{t-1}^{T})+m_{t}-\Pi^{-1}_{t}m_{t-1}.
\label{cons bc real}
\end{equation}
The steady state version of the consolidated budget constraint is the focus of this section. The existing real debt is denoted $b^{T}_{0}$, so the steady state net debt is $b=(1-\Omega)b_{0}^{T}$.
\begin{equation}
g-\tau+n(r^{n}-1)=(1-\Omega)b_{0}^{T}(r^{-1}-\Pi^{-1})+m(1-\Pi^{-1}).
\label{cons bc real ss}
\end{equation}
The following assumptions allow for interpretation of the constraint.
\begin{assumption}
Policymakers act to maintain the existing steady state real debt, so $b_{0}^{T}$ is fixed.
\end{assumption}
\begin{assumption}
The real rate given by $R=r\Pi^{-1}$ is fixed\footnote{This is a ``monetarist" assumption in \textcite{sargent1981unpleasant}}.
\end{assumption}
\begin{assumption}
The steady state demand for real currency $m=L(\Pi)$ is decreasing in inflation.
\end{assumption}
\begin{assumption}
The steady state quantity of real reserves $n=J(\Pi)$ is decreasing in inflation. 
\end{assumption}

Using the above and the steady state version of the interest rate rule (\ref{ior}), it is convenient to express the steady state budget constraint, showing the relationships to inflation as follows.
\begin{equation}
g-\tau+R^{b}(\Pi)+R^{n}(\Pi)=s^{m}(\Pi)
\label{cons bc real ss fig}
\end{equation}
The real interest payment on the net debt is
\[R^{b}(\Pi)=(1-\Omega)b_{0}^{T}(1-R^{-1})\Pi^{-1},\]
the real interest payment on reserves is
\[R^{n}(\Pi)=J(\Pi)\left[\left(\Pi R\right)^{\alpha_{i}}-1\right]\]
and the monetary seigniorage is
\[s^{m}(\Pi)=L(\Pi)(1-\Pi^{-1}).\]

The monetary seigniorage must cover the costs of the fiscal deficit, interest on existing debt and interest on reserves.

\begin{figure}[ht]
    \centering
    \includegraphics[width=10cm]{Fig 1.png}
    \caption{Representation of the steady state relation (\ref{cons bc real ss fig}) as functions of inflation $\Pi$.}
    \label{Fig1}
\end{figure}
The two sides of the constraint (\ref{cons bc real ss fig}) are represented in Figure \ref{Fig1}. The interest on the debt term $R^{b}(\Pi)$ is unambiguously non-increasing in inflation and the steady state deficit is constant, but the other two terms are more involved. The money demand $L(\Pi)$ function is decreasing with inflation, but the positive slope of the term $1-\Pi^{-1}$ dominates for $\Pi$ near one. As inflation rises the slope of the term $1-\Pi^{-1}$ approaches zero, so there is a maximum level of seigniorage. Whether the seigniorage function $s^{m}(\Pi)$ ultimately declines depends on the money demand $L(\Pi)$. In \textcite{sargent1981unpleasant} and \textcite{werning2024uma}, money demand is linear, which leads to a downward sloping portion of $s^{m}(\Pi)$ for large $\Pi$, but that is not guaranteed here\footnote{See \textcite{waters2024potentiallypleasant} for a related example where $s^{m}(\Pi)$ has a maximum but does not decline.}.

While the real interest on the debt is decreasing with inflation, the slope of the real interest on reserves $R^{n}$ is ambiguous. It is safe to assume that the slope of the payments on the left hand side of the steady state budget relation is less than that for seigniorage $s^{m}(\Pi)$ at the intersection for inflation near one.

\subsection{Results}

The steady state budget constraint (\ref{cons bc real ss fig}) represented in Figure \ref{Fig1} extends some established results to a more general setting.
\begin{result}
Higher deficits $g-\tau$ are associated with higher inflation, as shown by the $R^{b}+R^{n}+g-\tau$ curve shifting up, leading to a higher $\Pi^{*}$, but the curve has a maximum, so there is a limit to sustainable deficits.
\end{result}

MMT advocates have suggested that higher deficits do not necessarily lead to higher interest rates, which contradicts the short run intuition from a loanable funds graph. Figure 1 shows that the claim is not correct for steady states either. Given a fixed real rate, nominal rates and inflation both rise with the deficit.

\begin{result}
    \textbf{The Unpleasant Monetarist Arithmetic}:  \textcite{sargent2018recursive} describe ``one big open market operation" where the monetary authority sells bonds, raising the steady state net debt. On the graph, the fraction $\Omega$ falls, net debt $b$ rises, and the $R^{b}+R^{n}+g-\tau$ curve shifts up, leading to a higher $\Pi^{*}$.
\end{result}
\begin{result}
    \textbf{The Potentially Pleasant Modern Monetary Arithmetic}: Conversely, if the monetary authority commits to buying a fraction $\Omega$ of bonds, the $R^{b}+g-\tau$ curve shifts down, and inflation falls. Alternatively, for a given steady state inflation level $\Pi$, a higher fraction $\Omega$ of bonds purchased allows for a higher deficit $g-\tau$ (\cite{waters2024potentiallypleasant}). 
\end{result}
\begin{result}
    The greater the interest payments on the debt, the greater the benefit to monetary finance. If the term $R^{b}$ is small, a change in the parameter $\Omega$ has a modest impact. A low payment on the debt could arise due to a low level of the existing debt $b_{0}^{T}$ or a low steady state real rate $R$.
\end{result}
Secular stagnation, the long run fall in real rates, may be eroding the rationale for monetary finance.

The inclusion of interest on reserves in the constraint (\ref{cons bc real ss fig}) allows for some further insights.
\begin{result}
    A higher value of $\alpha_{i}$ in the policy rule (\ref{ior}) raises the cost of interest on reserves, shifts the costs curve up, and leads to higher steady state inflation $\Pi$.
\end{result}
\begin{result}
    For a given inflation target $\Pi$ and fraction of bonds purchased by the monetary authority $\Omega$, there is a trade-off between the costs to the government:  the fiscal deficit $g-\tau$ and the interest on reserves determined by $\alpha_{i}$. The higher the steady state fiscal deficit, the lower the interest on reserves must be and vice versa.
\end{result}

It is important to note that these are all steady state results, and short run intuitions could be different. In the short run, a monetary authority that buys bonds or lowers the policy rate can easily stimulate inflation due to liquidity effects. The fiscal authority might deviate from the steady state deficit but compensate in the future, a Ricardian policy. These points highlight the importance of credibility. For monetary finance to be beneficial, institutions must be sufficiently trusted for steady state values to be meaningful. Japan could be interpreted as such a country, but cases where monetary finance led to hyperinflation, such as those discussed in \textcite{sargent1982ends}, are cautionary tales. 

\subsection{On the monetary policy rate}

 Most countries claim to have a policy of monetary dominance, where fiscal authorities control deficits so that central banks can target inflation independently. Typically, monetary policy is described by a Taylor rules that includes a steady state natural rate of interest \textit{r-star} determined by the choice of inflation target, as in \textcite{Platanov2024flexpricedeterminacy}. In the present model, the steady state rate of inflation $\Pi$ and rate for interest on reserves $r^{n}$ are jointly determined by the fraction of bonds purchased by the monetary authority $\Omega$, the steady state deficit and the steady state rate of inflation according to the constraint (\ref{cons bc real ss}), along with the parameter $\alpha_{i}$. 
 
 The relationship between the steady state rate on reserves $r^{n}$ and market rate $r$ is governed by the policy parameter $\alpha_{i}$ according to the steady state version of the interest rate rule (\ref{ior}).
\begin{equation}
  r^{n}=r^{\alpha_{i}} \label{ior ss}  
\end{equation}

The choices for the parameters $\Omega$ and $\alpha_{i}$ by the monetary policymaker impact the cost of the interest on the debt and the interest on reserves and, therefore, the steady state inflation rate. To achieve monetary dominance the fiscal authority would have to set the steady state deficit $g-\tau$ according those choices and the desired inflation target $\Pi$.

The interest rate rule (\ref{ior}) and the steady state government budget constraint (\ref{cons bc real ss fig}) demonstrate the difference in short and long run dynamics with respect to inflation. The interest rate rate rule shows the short term conduct of policy where the policy would be raised in an attempt to reduce inflation above its steady state for $b_{\pi}>0$, which is necessary for stability, see Section 4.3. However, the steady state policy and inflation rates can move together due to exogenous shifts. Hence, an increase in the parameter $\alpha_{i}$ would increase the steady state inflation rate, but the resulting increase in the policy rate could lead to a short run decrease in inflation.

\section{The Model}

The next two sections develop a micro-founded model to demonstrate that the budget constraint (\ref{cons bc real}) and Assumptions 1-4 are consistent with household optimization. The model is a simplified version of a shopping time monetary model with a banking sector and interest on reserves as described in \textcite{ireland2014macroeconomics} that also allows for fiscal deficits. The model is simplified in that there is only one stochastic shock, and liquidity services for the household are given solely by bank deposits.

\subsection{Household Problem}

The household utility is determined by consumption and the disutility from
working and shopping time, which is alleviated by the liquidity service provided by
bank deposits $D_{t}$. \textcite{ireland2014macroeconomics} includes a cash infusion each period that also provides liquidity services, but all transactions are electronic here. The household does hold currency as one source for deposits.

The household enters time $t$ with currency held from
the previous period $M_{t-1}$, bank loans taken $L_{t}$ and the return from the bond holdings $B_{t-1}$ with yield $r_{t}$, so%
\begin{equation}
D_{t}=M_{t-1}+L_{t}+B_{t-1}-\frac{B_{t}}{r_{t}}.  \label{HHBC1}
\end{equation}%
In period $t$, the household pays gross interest $r_{t}^{l}$ on the loans and
receives gross interest $r_{t}^{d}$ on the deposits. Real consumption, the
price level, wage, hours worked and lump sum taxes are $c_{t}$, $P_{t}$, $%
W_{t},$ $h_{t},$ $\tau _{t}$. The household receives bank profits $E_{t}$ and firm profits $F_{t}$, though they do not enter the optimality decisions.
Hence, the household decision must also satisfy%
\begin{equation}
P_{t}c_{t}=W_{t}h_{t}+r_{t}^{d}D_{t}-r_{t}^{l}L_{t}-M_{t}+F_{t}+E_{t}-P_{t}%
\tau _{t}.  \label{HHBC2}
\end{equation}

The household maximizes%
\begin{equation}
E_{0}\sum\limits_{t=0}^{\infty } \beta ^{t} A_{t}\left[ \log{c}_{t}-\eta
\left( \left(1+\psi\right)^{-1}h_{t}^{\left(1+\psi\right)}+h_{t}^{s}\right) \right] ,  \label{U}
\end{equation}%
where the discount factor is $\beta <1$, and the preference shock $%
A_{t}=\exp a_{t}$ is such that%
\begin{equation}
a_{t}=\rho _{a}a_{t-1}+\varepsilon _{a,t} \text{ for }\varepsilon _{a,t}\sim
iid\left( 0\right).
\label{a}
\end{equation}
Shopping time is given by 
\begin{equation}
h_{t}^{s}=\chi ^{-1}\left( \frac{v_{a}P_{t}c_{t}}{D_{t}}\right) ^{\chi }.
\label{hs}
\end{equation}%
The parameter $v_{a}$ is positive and differentiates the disutility from
working and disutility from shopping. Again, all transactions are electronic so liquidity services
are provided by deposits, which reduces shopping time for $\chi>0$.

Writing the budget constraints (\ref{HHBC1}) and (\ref{HHBC2}), in real
terms using the convention $x_{t}=\frac{X_{t}}{P_{t}}$ yields%
\begin{eqnarray}
d_{t} &=&m_{t-1}\Pi_{t}^{-1}+l_{t}+b_{t-1}\Pi_{t}^{-1}-b_{t}r_{t}^{-1}  \label{realBC1} \\
c_{t} &=&w_{t}h_{t}+r_{t}^{d}d_{t}-r_{t}^{l}l_{t}-m_{t}+\frac{F_{t}}{P_{t}}+
\frac{E_{t}}{P_{t}}-\tau _{t}, \label{realBC2}
\end{eqnarray}%
where gross inflation is $\Pi _{t}$ as above. 

The objective (\ref{U}) at is maximized substituting for $h_{t}^{s}$
using equation (\ref{hs}), subject to the budget constraints (\ref{realBC1}) and (\ref{realBC2}), with Lagrange
multipliers $\lambda _{1,t}$ and $\lambda _{2,t}$. The following are the first order conditions with respect to
$c_{t},h_{t},d_{t},l_{t},b_{t},m_{t}$.%
\begin{eqnarray}
\beta^{t}A_{t}\left(1-\eta \chi h_{t}^{s}\right) &=&\lambda _{2,t}c_{t} \label{foc c} \\
\beta^{t}A_{t}\eta h_{t}^{\psi} &=&\lambda _{2,t}w_{t} \label{foc h} \\
\beta^{t}A_{t}\eta \chi h_{t}^{s} &=&d_{t}\left(\lambda _{1,t}-\lambda
_{2,t}r_{t}^{d}\right) \label{foc d} \\
\lambda _{1,t} &=&\lambda _{2,t}r_{t}^{l} \label{foc l} \\
\lambda _{1,t} &=& r_{t}E_{t}\left(\frac{\lambda _{1,t+1}}{\Pi_{t+1}}\right) \label{foc b}\\
\lambda _{2,t} &=&E_{t}\left(\frac{\lambda _{1,t+1}}{\Pi _{t+1}}\right) \label{foc m}
\end{eqnarray}

An immediate consequence of conditions (\ref{foc l}), (\ref{foc b}) and (\ref{foc m}) is the arbitrage condition that the rate on bonds and the rate on loans must be the same.
\begin{equation}
r^{l}_{t}=r_{t} \label{arb r}
\end{equation}

\subsection{Firm decision}

The production sector is standard and includes a competitive final goods
producing firms and monopolistically competitive intermediate goods
producing firms with quadratic price adjustment cost a l\`a  Rotemberg \cite{rotemberg1982stickyprices}.
Final goods $y_{t}$ are produced with constant returns to scale combinations
of intermediate goods $y_{t}\left( j\right) $ with price $P_{t}\left(
j\right) $, where the parameter $\theta $ represents the degree of
complementarity.%
\[
y_{t}=\left(\int\nolimits_{0}^{1}y_{t}\left( j\right)^{\dfrac{\theta-1}{%
\theta}}dj\right)^{\dfrac{\theta}{\theta-1}} 
\]%
Profits for the final goods firm is therefore%
\[
P_{t}\left( \int\nolimits_{0}^{1}y_{t}\left( j\right) ^{\dfrac{\theta -1}{%
\theta }}dj\right) ^{\dfrac{\theta }{\theta -1}}-\int\nolimits_{0}^{1}P_{t}%
\left( j\right) y_{t}\left( j\right) dj, 
\]%
so the demand for an intermediate good is given by%
\[
y_{t}\left( j\right) =\left( \frac{P_{t}\left( j\right) }{P_{t}}\right)
^{-\theta }y_{t}, 
\]%
and the final goods price level is
\[
P_{t}=\left( \int\nolimits_{0}^{1}P_{t}\left( j\right) ^{1-\theta
}dj\right) ^{\dfrac{1}{1-\theta }}.
\]
Intermediate goods production is determined by labor $h_{t}^{g}$ and constant productivity $z$.
\begin{equation}
y_{t}=zh_{t}^{g} \label{prod} \\
\end{equation}%

They also face the cost of adjusting prices%
\[
caj_{t}=\frac{\phi _{p}}{2}\left( \frac{P_{t}\left( j\right) }{\Pi
P_{t-1}\left( j\right) }-1\right) ^{2} y_{t}, 
\]%
where the parameter $\phi _{p}>0$ is positive and the steady state gross inflation
rate is $\Pi $.

Using the above relations, the time $t$ real profit for a representative
intermediate goods firm is%
\[
\frac{F_{t}\left( j\right) }{P_{t}}=\left( \frac{P_{t}\left( j\right) }{P_{t}%
}\right) ^{1-\theta }y_{t}-\left( \frac{P_{t}\left( j\right) }{P_{t}}\right)
^{-\theta }\frac{y_{t}}{Z_{t}}\frac{W_{t}}{P_{t}}-\frac{\phi _{p}}{2}\left( 
\frac{P_{t}\left( j\right) }{\Pi P_{t-1}\left( j\right) }%
-1\right) ^{2}y_{t}.
\]%
Maximizing the discounted sum of future profits%
\[
E_{0}\sum\limits_{t=0}^{\infty }\beta ^{t}\frac{F_{t}\left( j\right) }{P_{t}%
},
\]%
the first order condition with respect to $P_{t}\left( j\right) $ gives the
following.%


\begin{equation} \label{PC}
\begin{split}
0 &=\left( 1-\theta \right) \left( \frac{P_{t}\left( j\right) }{P_{t}}%
\right) ^{-\theta }\frac{y_{t}}{P_{t}}+\theta \left( \frac{P_{t}\left(
j\right) }{P_{t}}\right) ^{-\theta -1}\frac{y_{t}}{Z_{t}P_{t}}\frac{W_{t}}{%
P_{t}} \\
&-\phi _{p}\left( \frac{P_{t}\left( j\right) }{\Pi
P_{t-1}\left( j\right) }-1\right) \frac{y_{t}}{\Pi
P_{t}\left( j\right) } \\
&+\beta \phi _{p}E_{t}\left[ \left( \frac{P_{t+1}\left( j\right) }{\Pi P_{t}\left( j\right) }-1\right) \left( \frac{P_{t+1}\left(
j\right) }{\Pi P_{t}^{2}\left( j\right) }\right) y_{t+1}%
\right] .
\end{split}
\end{equation}

For a symmetric equilibrium where $P_{t}(j)=P_{t}$ for all $j$, the real profit is
\begin{equation} \label{real profit sym}
\frac{F_{t}}{P_{t}}=\left(1-w_{t}z^{-1}\right)y_{t}-caj(\Pi_{t})y_{t}
\end{equation}
where the real values of the wage and output are $w_{t},y_{t}$, and the cost of adjustment is a function of inflation.
\[
caj(\Pi_{t})=\frac{\phi_{p}}{2}\left(\frac{\Pi_{t}-\Pi}{\Pi_{t}}\right)^{2}
\]

The symmetric equilibrium for the optimality condition (\ref{PC}) is
\begin{equation} \label{PCsym}
\begin{split}
0 &=1-\theta+\theta(w_{t}Z_{t}^{-1})-\frac{\phi_{p}}{2}\left(\frac{\Pi_{t}-\Pi}{\Pi}\right)\left(\frac{\Pi_{t}}{\Pi}\right) \\
&+\beta\phi_{p}E_{t}\left[\frac{y_{t+1}}{y_{t}}\left(\frac{\Pi_{t+1}-\Pi}{\Pi}\right)\left(\frac{\Pi_{t+1}}
{\Pi}\right)\right]
\end{split}
\end{equation}
A log-linearized version has the form of a standard New Keynesian
Phillips Curve.

The resource constraint for the economy is
\begin{equation} \label{ycg}
    y_{t}=c_{t}+g_{t}+caj(\Pi_{t})y_{t}.
\end{equation}

\subsection{Banking Sector}

Banks take deposits and make loans while holding reserves. The technology
for generating both deposits and reserves requires labor input following
\textcite{christiano1995inside} and \textcite{ireland2014macroeconomics}.

Servicing real deposits $d_{t}=\dfrac{D_{t}}{P_{t}}$ requires both real reserves $n_{t}=\dfrac{N_{t}}{P_{t}}$ and
labor $h_{t}^{d}$ according to
\begin{equation} \label{d}
x_{a}\left( x_{n}^{\frac{1}{\nu }}n_{t}^{\frac{\nu -1}{\nu }}+\left(
1-x_{n}\right) ^{\frac{1}{\nu }}h_{t}^{d^{\frac{\nu -1}{\nu }}}\right) ^{%
\frac{\nu }{\nu -1}}\geq d_{t}.
\end{equation}

The parameter $\nu $ is the elasticity of substitution between the two
inputs and the parameter $x_{n}$ measures their relative importance. 

Servicing reserves also requires labor $h_{t}^{n}$ such that for a positive $%
\phi _{n}>0,$%
\begin{equation} \label{hn}
h_{t}^{n}\geq \phi _{n}n_{t}.
\end{equation}
Banks maximize real profits where they generate returns from loans and reserves and have costs
from wages and returns to depositors.
\begin{equation} \label{bank profits}
\frac{E_{t}}{P_{t}}=\left( r_{t}^{l}-1\right) l_{t}+\left(
r_{t}^{n}-1\right) n_{t}-\left( r_{t}^{d}-1\right) d_{t}-w_{t}h_{t}^{b}
\end{equation}%
The total labor employed by the bank $h_{t}^{b}$ is the sum of those
employed to service deposits and reserves. 
\begin{equation} \label{hb}
h_{t}^{b}=h_{t}^{d}+h_{t}^{n}
\end{equation}
The timing in the profits relation (\ref{bank profits}) is the same as with the household budget constraints (\ref{HHBC1}) and (\ref{HHBC2}) and the description of reserves after the budget constraint for the monetary authority (\ref{CB BC}).

Total labor supplied $h_{t}$ is the sum of that provided for goods production and banking.
\begin{equation} \label{h}
h_{t}=h_{t}^{g}+h_{t}^{b}
\end{equation}
Bank assets equal liabilities so
\begin{equation} \label{dnl}
d_{t}=n_{t}+l_{t}.
\end{equation}

By substituting for $l_{t}$ and $d_{t}$ with equations (\ref{d}), (\ref{dnl}) and (\ref{hn}), bank profits (\ref{bank profits}) can be written in term of $n_{t}$ and $h_{t}^{d}$.
\begin{equation}
\begin{split}
\frac{E_{t}}{P_{t}}& =(r_{t}^l-r_{t}^d)x_{a}\left[x_{n}^{1/\nu}n_{t}^{\frac{\nu-1}{\nu}}+(1-x_{n})^{1/\nu}(h_{t}^{d})^{\frac{\nu-1}{\nu}}\right]^{\frac{\nu}{\nu-1}} \\
&-(r_{t}^{l}-r_{t}^{n}+\phi_{n}w_{t})n_{t}-w_{t}h_{t}^{d}
\end{split}
\end{equation}

Maximizing bank profits, the first order conditions with respect to $n_{t}$ and $h_{t}^{d}$ follow.
\begin{eqnarray}
  n_{t}=(r^l_{t}-r^d_{t})^{\nu}x_{a}^{\nu-1}x_{n}d_{t}(r^{l}_{t}-r^{n}_{t}+\phi_{n}w_{t})^{-\nu}  \label{foc n} \\ 
  h^{d}_{t}=(r^l_{t}-r^d_{t})^{\nu}x_{a}^{\nu-1}(1-x_{n})d_{t}w_{t}^{-\nu},  \label{foc hd}
\end{eqnarray}

The quantities $n_{t}$ and $h_{t}^{d}$ must be positive, which implies that the interest rate on loans must be greater than that on deposits, $r^l>r^d$, matching the usual intuition for the way banks make profits. See \textcite{ireland2014macroeconomics} for a detailed discussion of the the optimal choices by the bank.

\subsection{Equilibrium Household Budget Constraints}

Given the optimizing behavior for households, firms and banks, the two household budget constraints can be written in terms of currency $m_{t}$ and reserves $n_{t}$. The combination of those two is equivalent to the consolidated government budget constraint (\ref{cons bc real}).

In equilibrium, the real quantities $b_{t}$, $m_{t}$, and $n_{t}$, are equal to the corresponding quantities in the budget constraints in Section 2. Substituting for real bank profits (\ref{bank profits}) in the household budget constraint (\ref{realBC2}) and firm profits (\ref{real profit sym}), and labor is split between goods production and banking (\ref{h}) yields
\[c_{t}=w_{t}h_{t}^{g}+d_{t}-l_{t}-m_{t}+\left(
r_{t}^{n}-1\right) n_{t}-\tau_{t}.\]
Using the production function (\ref{prod}), the resource constraint (\ref{ycg}) and the bank balance sheet condition (\ref{dnl}), the above constraint becomes
\begin{equation}
    m_{t}=g_{t}-\tau_{t}+r^{n}_{t}n_{t}.
    \label{HH BC m}
\end{equation}
Similarly, the first household budget constraint (\ref{realBC1}) with the bank condition (\ref{dnl}) can be written
\begin{equation}
    n_{t}=m_{t-1}\Pi_{t}^{-1}+b_{t-1}\Pi_{t}^{-1}-b_{t}r_{t}^{-1}.
    \label{HH BC n}
\end{equation}
Adding the above two equations (\ref{HH BC m}) and (\ref{HH BC n}) yields,
\begin{equation} \label{policy BC}
m_{t}-\frac{m_{t-1}}{\Pi_{t}}=g_{t}-\tau_{t}+\frac{b_{t-1}}{\Pi_t}-\frac{b_{t}}{r_{t}}+\left(r_{t}^{n}-1\right) n_{t},
\end{equation}
which is equivalent to the consolidated budget constraint (\ref{cons bc real}).

\subsection{The Euler equation}

An Euler equation is derived from the dynamic model, and the steady state version is consistent with Assumption 2. Solving out the Lagrange multipliers from the household optimization equations (\ref{foc h}), (\ref{foc b}), (\ref{foc l}), and using the interest rate arbitrage relationship (\ref{arb r}) gives the following.

\begin{equation}
    \beta E_{t}\left(\frac{r_{t+1}}{\Pi_{t+1}}\left(\frac{h_{t+1}}{h_{t}}\right)^{\psi}\frac{w_{t}}{w_{t+1}}\frac{A_{t+1}}{A_{t}}\right)=1 
    \label{Euler}
\end{equation}
The steady state version is
\begin{equation}
    \beta r\Pi^{-1}=1, \label{Euler ss}
\end{equation}
so the steady state real interest rate depends only on the constant discount factor $R=\beta^{-1}$, satisfying Assumption 2.

\subsection{Policy}

The policy specification allows for active monetary and fiscal policy. The fiscal authority responds to the monetary authority in the sense that bonds are issued in time $t$ without knowledge of the currency supply $M_{t}$. Government spending is exogenous and fixed, $g_{t}=g$. Taxes $\tau_{t}$ respond to push the real debt $b_{t}$ toward its steady state $b=(1-\Omega)b_{0}^{T}$ as in \textcite{eusepi2018fiscal}. The parameter $\phi_{b}$ represents how aggressiveness of the fiscal authority.
\begin{equation}
    \tau_{t}=\tau\left(\frac{b_{t}}{(1-\Omega)b_{0}^{T}}\right)^{\phi_{b}}
    \label{tau}
\end{equation}
The fiscal authority issues bonds according to real versions of the budget constraints (\ref{fiscal BC}) and (\ref{CB BC}), but its estimate of the transfer from the monetary authority $RCB_{t}$ uses lagged information about currency such that $E_{t-1}m_{t}=m_{t-1}$. Hence, the real net debt issued is given by the following.
\begin{equation}
    b_{t}=r_{t}\left[b_{t-1}\Pi_{t}^{-1}+g-\tau_{t}-m_{t-1}(1-\Pi_{t})+(r^{n}_{t}-1)n_{t}\right]
    \label{b issue}
\end{equation}

Monetary policy is determined by the interest rate rule (\ref{ior}). The monetary authority issues real currency $m_{t}$ to meet the demand given by the household condition (\ref{HH BC m}).

\section{Quantitative Results}

Calibration of the model allows for analysis of determinacy and expectational stability. In addition, this section shows quantitative counterparts of the steady state relations described in Section 2.

\subsection{Model Summary}

The model is simplified to describe the evolution of the 18x1 vector 
\[X_{t}=(y_{t},c_{t},d_{t},n_{t},\Pi_{t},r_{t},r^{d}_{t},r^{n}_{t},h_{t},h^{g}_{t},h^{d}_{t},h^{n}_{t},h^{s}_{t},w_{t},\tau_{t},b_{t},m_{t},a_{t})'. 
\]
Throughout the following specification, the bank balance sheet condition (\ref{dnl}) is used to substitute for $l_{t}$ and the arbitrage condition $r_{t}=r^{l}_{t}$ from (\ref{arb r}) is used to substitute for the interest rate on loans $r^{l}_{t}$.

The model equations include the specification for shopping time (\ref{hs}), the household optimality condition determined by equations (\ref{foc c}) and (\ref{foc h})
\begin{equation*}
    \chi h^{s}_{t}+\frac{c_{t}}{w_{t}}h_{t}^{\psi}=\eta^{-1},
\end{equation*}
the condition determined by equations (\ref{foc c}) and (\ref{foc h})
\begin{equation*}
    (r_{t}-r^{d}_{t})h_{t}^{\psi}d_{t}=  \chi w_{t}h^{s}_{t},
\end{equation*}
and the Euler equation (\ref{Euler}).
The model also includes the production function (\ref{production}), the Phillips Curve (\ref{PC}), the resource constraint (\ref{ycg}), as well as the banking conditions (\ref{d}), (\ref{foc n}), (\ref{foc hd}), (\ref{hn}) and the labor identity from combining (\ref{hb}) and (\ref{h}).
\begin{equation*}
    h_{t}=h^{g}_{t}+h^{d}_{t}+h^{n}_{t}
\end{equation*}
The model is closed with the budget conditions (\ref{HH BC m}), (\ref{HH BC n}), the interest rate rule (\ref{ior}), and the specification for the demand shock (\ref{a}). Fiscal policy is given by equations (\ref{tau}) and (\ref{b issue}).

\subsection{Calibration}

The model is calibrated to check for determinacy and expectational stability and examine steady state relationships. A number of the parameter values are taken from existing work. The household objective parameter values $\chi=1.0$ and $\psi=1.3$ are taken from the DSGE model of the Federal Reserve Bank of New York, \textcite{frbny_dsge_model}. Some banking sector values $\nu=0.25$ and $\phi_{n}=0.0005$ are from \textcite{ireland2014macroeconomics} as is the Phillips Curve parameter $\phi_{p}=50$. For the initial calibration, the value $\Omega=0.0$ reflects current policy, since most central banks do not systematically purchase bonds.

Steady state output $y$ and total hours $h+h^{s}$ are normalized to one. The banking parameter value $x_{n}=0.95$ is set higher than in \textcite{ireland2014macroeconomics} to achieve an value for labor hours in deposit creation $h^{d}$ in line with the other labor hours steady state values. The rest of the parameters are chosen to achieve reasonable steady state annual values: $\Pi=1.04$, $r=1.07$, $r^{d}=1.045$ and $r^{n}=1.02$. The values are $\beta=0.972$, $v_{a}=1.138$, $\eta=1.674$, $\nu=0.25$, $x_{a}=3.541$. The persistence of the shock given by $\rho_{a}=0.9$ is similar to that in \textcite{ireland2014macroeconomics}.

Since steady state output $y$ is one, the chosen values for steady state government spending $g=0.2$, taxes $\tau=0.18$ and fiscal deficit $g-\tau=0.02$ are interpreted as fractions of GDP. 

\subsection{Determinacy and Expectational stability}

The calibrated model is simulated with Dynare (\cite{adjemian2011dynare}) to verify steady state values and determine the solution under rational expectations. For all the parameter values studied, determinacy implies expectational stability. A model is determinate if its satisfies the conditions in  \textcite{blanchard1980solution}. The conditions for expectational stability are given by Proposition 10.3 in \textcite{evans2001learning}, see Appendix.

The monetary response to inflation is parameterized by $b_{\pi}$ in the interest rate rule (\ref{ior}), while the fiscal policy response to the debt is parameterized by $\phi_{b}$ in the tax rule (\ref{tau}).

 For passive fiscal policy $\phi_{b}=0$, expectational stability is achieved for $b_{\pi}>0.934$, similar to the Taylor principle-like results in \textcite{Platanov2024flexpricedeterminacy}, \textcite{bullardmitra02} and \textcite{eusepi2018fiscal}. A pegged policy rate $r^{n}_{t}=r^{n}$, when $b_{\pi}=0$ does not achieve stability, following \textcite{friedman1968role}, \textcite{howitt1992interest}, and \textcite{evans2018interest}. For an alternative parameterization\footnote{The new values are $\alpha_{i} = 0.292$, $\beta = 0.986$, $\eta = 1.672$, $v_{a} = 1.917$, $x_{a} = 4.828$.} where $\Pi=1.02$, the results are similar and the threshold for stability is $b_{\pi}>0.96$ when $\phi_{b}=0$.

Active fiscal policy plays a minimal role in achieving stability. The parameter $\phi_{b}$ must be positive to stay in an active monetary/passive fiscal policy regime, in the language of \textcite{leeper1991active}, but it must be lower than 0.153 to achieve stability. Overly active fiscal policy can be de-stabilizing, as in \textcite{eusepi2018fiscal}. Whether these results are robust in models with richer dynamics is a topic for future work. 

So a direct policy rate response to inflation is required for expectational stability, which verifies the usual intuition for the short run response of the policy rate to inflation discussed in Section 2.4. Raising the policy rate acts to bring inflation down toward the target. For the steady state analysis, inflation and interest rates move together in response to parameter shifts (except for $\beta=R^{-1}$).

\subsection{Steady State Relationships}

This section reports results about the relationships described in the steady state budget constraint (\ref{cons bc real ss fig}). The ultimate goal is to show quantitatively the trade-off between the steady state fiscal deficit $g-\tau$ and the monetary policy parameter $\alpha_{i}$, which determines the cost of the interest on reserves, for a given inflation target $\Pi$, as described in Result 6.

\subsubsection{Deficits versus Inflation}

Figure \ref{Fig 2} verifies the relationships between the deficit and inflation as in Result 1 and the relationship between inflation and money in Assumption 3. Values of the real tax $\tau$ range from 0.3 to 0.02, which translates into an increasing fiscal deficit $g-\tau$ on the horizontal axis. Inflation $\Pi$ rises with the deficit while currency $m$ falls. The analogous graphs for reserves is very similar, so the model is consistent with Assumptions 3 and 4. Steady state deposits are decreasing with inflation as well.

\begin{figure}[h]
\centering 
        \includegraphics[width=9cm]{Fig 2.png}    
    \caption{The upward sloping steady state inflation $\Pi$ and downward sloping steady state currency $m$ are functions of the fiscal deficit $g-\tau$. The net debt $b=(1-\Omega)b^{T}_{0}$ takes values according to $\Omega=0.0, 0.5, 1.0$ and $b^{T}_{0}=1.0$.}
       \label{Fig 2}
\end{figure}

The benefit of the monetary authority purchasing bonds is apparent. The net debt $(1-\Omega)b^{T}_{0}$ varies with the choice of $\Omega$. The existing debt is set to $b^{T}_{0}=1.0$, or 100\% of output, and the fraction purchased by the monetary authority is set to 0.0, 0.5 and 1.0. Higher values of $\Omega$ give lower net debt and lower steady state inflation rates, as in Result 3.

\renewcommand{\arraystretch}{1.75} 

\begin{table}[h]
\setlength{\tabcolsep}{3pt}
\begin{tabular}{lrccccccccccc}
\hline
                                                         & \multicolumn{1}{l}{}  & \multicolumn{11}{c}{\large{\textbf{$g-\tau$}}}                                                                                                                                                                                                                                                                                                                                                                                                                                                            \\ \hline
\multicolumn{1}{l}{}                                    & \multicolumn{1}{l}{}  & \multicolumn{1}{c}{\textit{\textbf{-0.04}}} & \multicolumn{1}{c}{\textit{\textbf{-0.03}}} & \multicolumn{1}{c}{\textit{\textbf{-0.02}}} & \multicolumn{1}{c}{\textit{\textbf{-0.01}}} & \multicolumn{1}{c}{\textit{\textbf{0.0}}} & \multicolumn{1}{c}{\textit{\textbf{0.01}}} & \multicolumn{1}{c}{\textit{\textbf{0.02}}} & \multicolumn{1}{c}{\textit{\textbf{0.03}}} & \multicolumn{1}{c}{\textit{\textbf{0.04}}} & \multicolumn{1}{c}{\textit{\textbf{0.05}}} & \textit{\textbf{0.06}} \\ \hline
\multicolumn{1}{l}{\multirow{5}{*}{\large{\textbf{$\Omega$}}}} & \textit{\textbf{0.0}}  & \multicolumn{1}{c}{1.07}                    & \multicolumn{1}{c}{1.49}                    & \multicolumn{1}{c}{1.93}                    & \multicolumn{1}{c}{2.40}                    & \multicolumn{1}{c}{2.91}                & \multicolumn{1}{c}{3.44}                   & \multicolumn{1}{c}{4.01}                   & \multicolumn{1}{c}{4.61}                   & \multicolumn{1}{c}{5.24}                   & \multicolumn{1}{c}{5.91}                   & 6.61                   \\ \cline{2-13} 
\multicolumn{1}{l}{}                                   & \textit{\textbf{0.25}} & \multicolumn{1}{c}{0.76}                    & \multicolumn{1}{c}{1.16}                    & \multicolumn{1}{c}{1.58}                    & \multicolumn{1}{c}{2.02}                    & \multicolumn{1}{c}{2.50}                & \multicolumn{1}{c}{3.01}                   & \multicolumn{1}{c}{3.55}                   & \multicolumn{1}{c}{4.13}                   & \multicolumn{1}{c}{4.73}                   & \multicolumn{1}{c}{5.37}                   & 6.05                   \\ \cline{2-13} 
\multicolumn{1}{l}{}                                   & \textit{\textbf{0.5}}  & \multicolumn{1}{c}{0.47}                    & \multicolumn{1}{c}{0.84}                    & \multicolumn{1}{c}{1.24}                    & \multicolumn{1}{c}{1.66}                    & \multicolumn{1}{c}{2.12}                & \multicolumn{1}{c}{2.60}                   & \multicolumn{1}{c}{3.12}                   & \multicolumn{1}{c}{3.67}                   & \multicolumn{1}{c}{4.25}                   & \multicolumn{1}{c}{4.86}                   & 5.51                   \\ \cline{2-13} 
\multicolumn{1}{l}{}                                   & \textit{\textbf{0.75}} & \multicolumn{1}{c}{0.19}                    & \multicolumn{1}{c}{0.54}                    & \multicolumn{1}{c}{0.92}                    & \multicolumn{1}{c}{1.32}                    & \multicolumn{1}{c}{1.75}                & \multicolumn{1}{c}{2.21}                   & \multicolumn{1}{c}{2.70}                   & \multicolumn{1}{c}{3.22}                   & \multicolumn{1}{c}{3.78}                   & \multicolumn{1}{c}{4.37}                   & 4.99                   \\ \cline{2-13} 
\multicolumn{1}{l}{}                                   & \textit{\textbf{1.0}}  & \multicolumn{1}{c}{-0.68}                   & \multicolumn{1}{c}{0.26}                    & \multicolumn{1}{c}{0.62}                    & \multicolumn{1}{c}{1.00}                    & \multicolumn{1}{c}{1.40}                & \multicolumn{1}{c}{1.84}                   & \multicolumn{1}{c}{2.31}                   & \multicolumn{1}{c}{2.81}                   & \multicolumn{1}{c}{3.33}                   & \multicolumn{1}{c}{3.89}                   & 4.49                   \\ \hline
\end{tabular}
\vspace{3mm}
\caption{Steady State inflation rates as simple annual percentages for given deficits $g-\tau$ and fractions $\Omega$ of debt purchased by the monetary authority. The parameter $\alpha_{i}=0.3$.}
\label{tab:inf}
\end{table}

\begin{figure}[h]
\centering
    \includegraphics[width=10cm]{Fig 53.png}
    \caption{Steady state seigniorage $s^{m}$, interest on debt $R^{b}$ and interest on reserves $R^{n}$ as a function of inflation $\Pi$, for taxes $\tau$ varying from 0.3 to 0.0. The fraction of bonds purchased by the monetary authority is set to zero, $\Omega=0.0$.}
    \label{Fig53}
\end{figure}

Table \ref{tab:inf} shows the annualized simple inflation rates $\Pi-1$ for varying levels of the fraction $\Omega$ and the deficit $g-\tau$. Starting from a deficit of 0.02, or 2\% of GDP, and no debt purchased by the monetary authority $\Omega=0.0$ where inflation is 4\%, the benefit of switching to an MMT policy $\Omega=1.0$ would be a 1.7\% drop in inflation. Alternatively, full monetary finance would allow for an increase in the deficit to over 5\% of GDP while maintaining a 4\% inflation rate.

\subsubsection{Seigniorage, interest on the debt and interest on reserves versus inflation}

Figure 1 in Section 2 represents the constraint (\ref{cons bc real ss fig}), which contains much of the intuition. Figure \ref{Fig53} shows steady state seigniorage, interest on the debt and interest on reserves as functions of inflation for the calibrated model. Again, the driving factor is a decrease in taxes, so the deficit $g-\tau$ is not included on the graph. One can think of the curves in Figure \ref{Fig53} as intersections with an increasing deficit line as in Figure 1.

Figure \ref{Fig53} shows the case without debt purchases by the monetary authority $\Omega=0.0$, so the interest on the net debt $R^{b}$ is positive. It is declining with inflation as shown by the specification in Section 2, but the slope is barely discernible on the graph. The slope of the interest on reserves curve $R^{n}$ is positive but less than the seigniorage curve $s^{m}$ near equilibrium inflation, as required. In the case under the MMT policy where the monetary authority buys all the bonds $\Omega=1.0$, the graph is very similar, but the interest on the debt $R^{b}$ is constant at zero, since the public holds none of the bonds.

\subsubsection{Targeting Inflation}

Steady state inflation is increasing with any parameter that directly affects a cost to the government, including the deficit $g-\tau$ and the monetary policy parameters $\alpha_{i}$, which determines the steady state interest rate on reserves $r^{n}$. Figure \ref{Fig56} shows the combinations $(g-\tau,\alpha_{i})$ that give steady state simple inflation levels of 2\% and 4\%, for both cases $\Omega=0.0,1.0$. There is a non-linear trade-off between the cost of the deficit and the cost of interest on reserves, reflecting Result 6. For a given inflation target, the higher the deficit, the lower the steady state interest on reserves must be and vice versa. Monetary finance, a higher $\Omega$, allows for higher costs for the same inflation target.

\begin{figure}[h]
    \includegraphics[width=10cm]{Fig 56.png}
    \caption{The plots are  combinations $(g-\tau,\alpha_{i})$ that give steady state simple inflation levels of 2\% and 4\% in the cases of full monetary finance and no monetary finance. The two dots represent the 2\% inflation case where $\alpha_{i}=0.19$, which corresponds to $r^{n}=1.01$.}
    \label{Fig56}
\end{figure}

A simple example shows the maximum benefit of monetary finance. If the inflation target is 2\% and a steady state rate on reserves of 1\% (simple, annual for both) translates to the value $\alpha_{i}=0.19$. Without any monetary finance $\Omega=0.0$, a balanced budget is required, $g-\tau=0.0$, represented by the large dot on the left in Figure \ref{Fig56}. For the same inflation target and policy rate, switching to full monetary finance, where $\Omega=1.0$, allows for a fiscal deficit $g-\tau=0.03$. So the switch to full monetary finance allows for an increase in the steady state deficit of 3\%, which is the saving from avoiding the interest on the debt, following Result 3.

The 3\% increase in the deficit made possible by full monetary finance is the real interest on the debt $R^{b}$ in equation \ref{cons bc real ss fig} and Figure \ref{Fig1}. If the fraction of bonds $\Omega$ changes from zero to one, the $R^{b}$ term is eliminated. To maintain the same steady state inflation rate $\Pi^{*}$, the deficit $g-\tau$ must rise by $R^{b}(\Pi)=b_{0}(1-R^{-1})\Pi^{*^{-1}}$.

\subsection{Discussion}

MMT advocates raise an important point about the benefits of monetary finance. However, there are a number of reasons for caution, particularly concerning their focus on full monetary finance, which would eliminate the debt held by the public. Partial monetary finance in combination with targeting the interest rate on reserves is potentially beneficial, particularly for countries with a high burden from the interest on their debt.

Monetary finance leads to lower inflation or allows higher sustainable government cost, but it must be emphasized that these results come from steady state analysis. Suggestions that the interest on the debt can be eliminated (\textcite{globerman2021mmt}) immediately are not correct. Decades of experience with monetary policy conducted through open market operations indicate that, if the Federal Reserve were to buy all the outstanding U.S. debt over a short period of time, higher inflation would almost certainly result to due to liquidity effects. A switch to monetary financing should not be done quickly.

The existing U.S. debt is the largest stock of safe assets and is an important part of the world financial market. Full monetary finance eliminates the bond market. MMT advocates say that the government should still issue bonds (\cite{mitchell2019mmt}), but there is little discussion about the quantity or the use of the funds raised. As \textcite{Woodford90debt}, \textcite{AIYAGARI98debt} and \textcite{brunnermeier2020bubble} show, government debt provides safe assets that help overcome financial market imperfections. A government is not required to pay off the existing debt and may have good reason not to. For example, Norway maintains a debt of over 30\%  of GDP though its sovereign wealth fund could buy all of it.

\textcite{AIYAGARI98debt} compute the optimal debt level to be about two-thirds output. Consider the U.S. values\footnote{Values come from the \textcite{USTreasury2025Interest}} in 2025:  an inflation target of 2\%, net debt of 0.966 of output, and an interest rate on the debt of 3.3\%., which implies a real rate of 1.27\%. \textbf{If partial monetary finance is used to reduce the net debt from 0.966 to 0.666 of output, the fraction purchased is $\Omega=0.31$. The savings on the interest on debt, and hence the increase in the steady state deficit that could be sustained for the U.S., is 0.32\% of output.}

For the case of Japan, though the gross debt ratio is over 250\% of GDP, their net debt was 94.3\% in 2023, similar to the U.S., see \textcite{chien2023japan}. More importantly, the real rate on the debt was negative so there is no incentive for further monetary finance.

\subsection{Extensions}

The DSGE model described here could provide the basis for empirical analysis of countries that engage in monetary finance, such as Japan and Argentina. 

Future work should extend the stability results to a model with price stickiness and interest rate rules that include a response to output and interest rate smoothing.

Reconciling the short run and long run impacts of policy for the present model is an important goal. For example, a permanent tax cut would increase the deficit putting upward pressure on inflation. The interest rate rule indicates that the monetary authority must raise rates to fight inflation while lowering the steady state rate on reserves to compensate for the cost to the government of the higher deficits. The proper path of the policy rate over the transition that accounts for both considerations is a non-trivial issue.

Another question the present work raises but cannot answer is the the optimal choice of $\alpha_{i}$. The model suggests that lower values reduce the cost of interest on reserves, allowing for higher deficits, and nothing rules out negative values of $\alpha_{i}$. In practice, the zero lower bound rules them out, and the results above indicate stability would be a problem when $r^{n}_{t}=0$. Hence, the inflation target and the steady state interest on reserve $r^{n}$ are assumed to be positive, so the value of $\alpha_{i}$ must be positive as well. A lower value reduces the cost of interest on reserves, but means the zero lower bound is more likely to bind. 

The values of $g-\tau, \alpha_{i}, \Omega, b_{0}^{T},\Pi$ are all interdependent policy choices. The ultimate goal is to develop an empirically robust model that reflects the welfare implications of these choices and compute the optimum values given the budget constraint (\ref{cons bc real ss fig}).

\section{Estimation}

This section reports results of Bayesian estimation of an extension of the model above using data from the U.S., 2009Q1-2025Q4, the sample where monetary policy has been conducted through the interest rate on reserves. Such estimations are inherently challenging due to the short sample that includes long stretches where the zero lower bound (ZLB) on interest rates binds and the severe recession due to COVID-19. Nevertheless, the data are consistent with the model and the estimations are informative, particularly regarding monetary policy parameters.

The model is extended to include additional stochastic and dynamic elements. The interest rate rule (\ref{ior}) for monetary policy includes a lower bound, shocks $z^r_t$ and responds to output.
\begin{equation}
  r^{n}_{t} = \max\left[1,r_{t}^{a_i}\left(\frac{\pi_t}{\pi}\right)^{b_\pi}\left(\frac{y_t}{y}\right)^{b_y} e^{z^r_t}\right]
  \label{iorest}
\end{equation}

The tax rule (\ref{tau}) is extended to included shocks and persistence parameterized by $\rho_\tau$.
\begin{equation}
   \tau_t = \tau_{t-1}^{\rho_\tau}\,\tau^{1-\rho_\tau}\left(\frac{b_{t-1}}{(1-\omega)\,b_0}\right)^{\!\varphi_b} e^{\varepsilon^\tau_t}
  \label{tauest}
\end{equation}

The extension of equation (\ref{b issue}) describing the evolution of government debt $b_{t}$ includes the endogenous tax variable and a stochastic government spending variable $g_{t}$. It also uses the assumption ratio $\pi_{t}/r_{t}=\beta$ is fixed at its steady state value so that the debt variable $b_{t}$ is predetermined.
\begin{equation}
 b_t = \frac{r}{\pi}b_{t-1}+ b_0\left(1 - \frac{r}{\pi}\right) + (g_{t} - g) - ( \tau_t -\tau)
  \label{b issue est}
\end{equation}

Besides the shock to preferences $A_{t}$ in household utility (\ref{U}) and the stochastic terms $\varepsilon^\tau_t$, $g_{t}$, the productivity parameter $z$ in (\ref{prod}), the money demand parameter $v_{a}$ in (\ref{hs}) and the deposit demand parameter $x_{a}$ in (\ref{d}) are replaced with stochastic variables. All the stochastic terms are AR(1) processes where the $\rho\in [0,1)$ parameters represent the persistence in each.
\begin{align*}
a_t &= a_{t-1}^{\rho_a} \exp(\varepsilon_{a,t}) \\
z_t &= z_{t-1}^{\rho_z} z^{(1-\rho_z)} \exp(\varepsilon_{z,t}) \\
v_{a,t} &= v_{a,t-1}^{\rho_{va}} v_{a}^{(1-\rho_{va})} \exp(\varepsilon_{v_{a,t}}) \\
x_{a,t} &= x_{a,t-1}^{\rho_{x_{a}}} x_{a}^{(1-\rho_{xa})} \exp(\varepsilon_{x_{a,t}}) \\
g_t &= \rho_g g_{t-1} + (1-\rho_g) g + \varepsilon_{g,t} \\
z_t^r &= b^r z_{t-1}^r + \varepsilon^r
\end{align*}

\subsection{Data}

The model is estimated on seven quarterly U.S.\ time series spanning 2009Q1--2025Q4 (68 observations).  All series are extracted from the Federal Reserve Bank of St.\ Louis FRED database.  Table~\ref{tab:data} in Appendix~\ref{app:data} provides a complete description of each series, including the FRED mnemonic, source agency, and transformation applied.

The seven observables entering the likelihood are:
\begin{enumerate}
  \item \textbf{Output} ($y^{obs}$): log of real GDP divided by nonfarm business sector labor productivity.
  \item \textbf{Deposits} ($d^{obs}$): log of deposits at commercial banks divided by labor productivity.
   \item \textbf{Currency} ($m^{obs}$): log of currency in circulation divided by labor productivity.
  \item \textbf{Inflation} ($\pi^{obs}$): log quarterly gross change in the GDP deflator.
  \item \textbf{Loan rate} ($r^{obs}$): log quarterly gross prime lending rate.
  \item \textbf{Policy rate} ($r^{n,obs}$): log quarterly gross effective federal funds rate.
  \item \textbf{Tax rate} ($\tau^{obs}$): log of federal tax receipts as a share of nominal GDP.
\end{enumerate}

The sample includes two distinct ZLB episodes: 2009Q1--2015Q4 (the post-Great Financial Crisis period) and 2020Q1--2021Q4 (the COVID-19 period).  The sample mean of the quarterly log gross inflation rate corresponds to approximately 2.2\% annualized inflation.

All series are demeaned by their full sample mean prior to estimation.  This ensures that the observables have zero unconditional mean, consistent with the model's measurement equations, which express each observable as a log-deviation from its steady-state value plus a measurement error. Output and deposits are scaled using labor productivity, which is standard.

\subsection{Calibration}

The calibrated parameters must be adjusted to account for the quarterly data. The discount factor $\beta$ is the most consequential calibration choice.  It determines the steady-state real interest rate through the Euler equation $r=\pi \beta$.

The ZLB is another factor in the choice of $\beta$. The estimation imposes the bound on the policy rate $r^{n}_{t}$ directly, but the deposit rate must be above the bound as well. The value $\beta = 0.9967$ is chosen so that the steady state bond/loan rate $r$ is above the inflation target $1.02^{0.25}$ but the steady state deposit rate is also above the ZLB, $r^{d}>1$.

Table \ref{tab:calib} in Appendix \ref{app:data} shows the values of the other calibrated parameters. Since achieving convergence is difficult for the fiscal feedback parameter, it is set to $\phi_{b}=0.05$, since a positive value is required to achieve determinacy in the extended model. Robustness tests using other values indicate that 0.05 is reasonable.

The initial estimation uses a steady state tax $\tau=0.18$, matching the calibration used for the steady state analysis. The results ( Table \ref{tab:posteriors_all_018} in Appendix \ref{app:tau018}) indicate that the implied steady state inflation is unrealistically high (approximately 5\%), so estimations for $\tau=0.20$ are given more attention, Table \ref{tab:posteriors_020} shows the prior and posterior means and distributions for the estimated parameters.

\subsection{Estimation}

The model is estimated using the Metropolis--Hastings algorithm implemented in Dynare~7 (\textcite{adjemian2011dynare}.  The posterior mode is located using the \texttt{csminwel} numerical optimizer.  The covariance of the jumping distribution is Initialized from the prior variance matrix.  Two parallel chains of 50,000 draws each are generated with a scaling factor of 0.20.  In total, 14 parameters are estimated: 7 structural parameters and 7 shock standard deviations.

The ZLB is imposed on the policy rate $r^{n}_{t}$ using the OccBin methodology described in \textcite{GuerrieriIacoviello2015}. Since the model at the ZLB is not guaranteed to be determinate, the estimation method uses the piecewise linear solution.

\begin{table}[hp]
\centering
\small
\begin{tabular}{lll@{\hspace{10pt}}c@{\hspace{10pt}}ccc}
\toprule
Symbol & Description & Distribution & Prior mean & Prior std & Post.\ mean & 90\% HPD \\
\midrule
$v^a$    & Preference scale                & Gamma  & 2.50 & 0.75 & 2.757 & $[2.10,\; 3.42]$ \\
$\eta$   & Labor preference               & Gamma  & 1.53 & 0.40 & 1.244 & $[0.84,\; 2.00]$ \\
$b_\pi$  & Inflation response     & Normal & 1.50 & 0.25 & 2.122 & $[1.82,\; 2.53]$ \\
$b_y$    & Output response       & Gamma  & 0.50 & 0.25 & 0.422 & $[0.17,\; 0.59]$ \\
$a_i$    & Policy-rate level weight        & Beta   & 0.50 & 0.15 & 0.503 & $[0.21,\; 0.64]$ \\
$b_r$    & MP shock persistence            & Beta   & 0.70 & 0.15 & 0.441 & $[0.28,\; 0.65]$ \\
$\rho_z$ & Productivity persistence  & Beta   & 0.50 & 0.15 & 0.401 & $[0.19,\; 0.71]$ \\
\bottomrule
\multicolumn{4}{l}{\parbox{0.82\linewidth}}
\end{tabular}
\caption{Posterior Estimates for 2 chains with 50,000 draws ($\varphi_b = 0.05$, $\tau = 0.20$)}
\label{tab:posteriors_020}
\end{table}

\begin{table}[hp]
\centering
\small
\begin{tabular}{llccc}
\toprule
Symbol & Description & Posterior mean & Geweke $p$ (15\% taper) & Converged? \\
\midrule

$b_\pi$  & Inflation response    & 2.122 & 0.002 & No  \\
$b_y$    & Output response       & 0.422 & 0.933 & Yes \\
$b_r$    & MP shock persistence           & 0.441 & 0.045 & No  \\
$a_i$    & Policy-rate level weight       & 0.503 & 0.411 & Yes \\
$v^a$    & Preference scale               & 2.757 & 0.505 & Yes \\
$\eta$   & Labor preference               & 1.244 & 0.003 & No  \\
$\rho_z$ & Productivity persistence & 0.401 & 0.034 & No  \\
\bottomrule
\end{tabular}
\caption{Geweke Convergence $p$-Values ($\varphi_b=0.05$, $\tau=0.20$)}
\label{tab:geweke}
\end{table}

\begin{figure}[htbp]
\centering
\includegraphics[width=0.95\textwidth]{posteriors_2x2.pdf}
\caption{Prior (dashed) and posterior (shaded/solid) densities for the preferred specification ($\varphi_b = 0.05$, $\tau = 0.20$).  Vertical lines mark posterior means.  The posterior for $b_\pi$ is concentrated well above unity.  50\,000 draws per chain $\times$ 2 chains.}
\label{fig:ppost}
\end{figure}

Table \ref{tab:posteriors_020} shows the posterior means and the Highest Posterior Density intervals (HPD). Figure \ref{fig:ppost} shows the prior and posterior densities for the interest rate rule parameters. Table \ref{tab:geweke} shows the Geweke diagnostic $p$-values for convergence as described in \textcite{Geweke1992}. Although four of the parameters do not converge, the policy rule parameter $b_{y}$ does, and the posterior distribution shows the HPD for $b_{\pi}$ is narrow so the estimation is informative. Furthermore, the estimated value of $b_{\pi}=2.12$ is robust in estimations with alternative calibrations. The densities show that estimation of the parameter $\alpha_{i}$ is not particularly informative, though it is close to the value derived from sample means of inflation and the policy rate.

\subsubsection{Alternative fiscal feedback parameters}

Table \ref{tab:posteriors_all_020} in Appendix \ref{app:tau020}, shows posterior means for different choices of the fiscal feedback parameter $\varphi_b$. The Taylor-rule parameters are stable:  under $\tau = 0.20$: $b_\pi \in [2.01,1.94,2.01]$, $b_y \in [0.35,0.29,0.40]$ and $b_r \in [0.40,0.38,0.39]$ for $\varphi_b \in \{0.05, 0.10, 0.15\}$. The Marginal Data Density (MDD) diagnostic statistics show better performance for $\varphi_b=0.10$, but that is an artifact of the near indeterminacy for that estimation, so the value $\varphi_b=0.05$ is still preferred. For the case where $\tau = 0.18$ the estimates of the policy parameters and the MDD statistics are even more stable. See Table \ref{tab:posteriors_all_018} in Appendix \ref{app:tau018}.

\subsection{Determinacy and Expectational Stability}

The results for determinacy and expectational stability are similar to those for the calibrated model in the previous section. Again, determinacy is given the conditions in \textcite{blanchard1980solution}, and the expectational stability condition are given by Proposition 10.3 in \textcite{evans2001learning}. Figure \ref{fig:det_estab_combined} shows results for the policy parameters $(b_\pi, b_y) \in [0,3]$ for the posterior means from the estimation with $(\tau,\phi_{b})=(0.20,0.05)$.
\begin{figure}[h!]
\centering
\includegraphics[width=0.85\textwidth]{det_estab_combined.pdf}
\caption{Classification of determinacy and expectational stability given estimated parameters for the case $\varphi_b = 0.05$, $\tau = 0.20$}
\label{fig:det_estab_combined}
\end{figure}

The regions of determinacy and expectational stability coincide and are completely determined by the response to inflation $b_{\pi}$. A response slightly stronger than indicated by the Taylor principle $b_{\pi}=1$ achieves stability. This required response is stronger than the calibrated case since the fiscal feedback parameter $\phi_{b}$ is positive here. Even in the face of multiple crises, U.S. policy appears to satisfy conditions for maintaining stability.

\subsection{Estimation discussion}

The analysis takes on the task of estimating a micro-founded DSGE model where monetary policy is conducted through the interest rate on reserves for a sample that includes two major recessions that led to two stretches at the ZLB and one spike in inflation. The model gives intuitive estimates, particularly about the parameters in the interest rate rule, which indicates that policy in practice is well within the boundaries to achieve stability.

The estimations do rely on a number of calibrated values, and there is scope for estimating more parameters in future work. Incorporating data on the deficit and debt is another potential avenue for future work. Using data on currency and reserve would also be desirable, though there are major difficulties to overcome. In particular, reserves have risen a fluctuated dramatically over the sample, partly due to the ZLB, but also due to institutional changes detailed in \textcite{Duffie2026}.

\section{Conclusion}

Partial monetary finance and stabilization policy conducted through the interest rate on bank reserves is a policy combination used in Japan 2012-2020 that deserves attention. Full monetization of government debt is not recommended, since the bond market has an important role in the financial system. However, for countries with debt in excess of that needed for financial markets, partial monetary finance could alleviate some of the cost of the interest on the debt. Monetary finance could interfere with open market operations but  not monetary policy conducted through the interest rate on bank reserves, which is current practice in many countries.

The steady state consolidated budget constraint is derived from a dynamic model that achieves determinacy and expectational stability for monetary policy rule where the interest rate on reserves satisfies the Taylor Principle. Bayesian estimation verifies that U.S. policy satisfies these conditions for the sample when monetary policy targets the interest rate on reserves. The steady state budget constraint shows that monetary seigniorage must cover the costs from the interest on the debt, the fiscal deficits and the interest on reserves. There is an inverse relationship between the fraction of bonds purchased by the monetary authority and the steady state inflation rate, which demonstrates the benefit of monetary finance and the unpleasant monetarist arithmetic. Alternatively, for a given inflation target, there is a policy trade-off between the fiscal deficit and the steady state interest rate on reserves.

That these are steady state results emphasizes the importance of credibility. They do not justify the monetary finance episodes that led to accelerating inflation in countries such as Argentina. In those cases, bond purchases by the monetary authorities increased with fiscal deficits and interest on the debt, so there were no meaningful steady state values for those costs. In contrast, the trust in Japanese institutions has allowed them to use monetary finance as a tool for stability and to contain the impact of sovereign debt. For countries without such trust, the partial monetary finance policy described here may require an independent fiscal authority with the power to control deficits.

\printbibliography

\clearpage
%===========================================================
\appendix
\section{Expectational Stability}
%===========================================================

 Expectational stability governs the stability of models for a wide range of non-rational methods of forming expectations. The log-linearized model has the form
\begin{equation}
    x_{t} = A_{0} + A_{1}E_{t-1}x_{t+1} + A_{2}x_{t-1},
    \label{model}
\end{equation}
where the vector $x_{t}$ contains the log-linearized variables from $X_{t}$, including the shock $a_{t}$ from (\ref{a}).
The minimum state variables rational expectations solution has the form
\begin{equation}
    x_{t}=B_{0}+B_{1}x_{t-1}.
    \label{sol}
\end{equation}
The results in \textcite{evans2001learning} (Chapter 10) determine the expectational stability of the model (\ref{model}). Define the matrices $DT_{a}$ and $DT_{b}$ as follows.
\begin{eqnarray}
     & DT_{a}=(I-A_{1}B_{1})^{-1}A_{1} \\
     & DT_{b}=[(I-A_{1}B_{1})^{-1}A_{2}]\otimes[(I-A_{1}B_{1})^{-1}A_{1}]
\end{eqnarray}

\begin{proposition}
  [\textcite{evans2001learning}, Proposition 10.3] Let the value $\Lambda$ be maximum of the real parts of the eigenvalues of $DT_{a}$ and $DT_{b}$. The model (\ref{model}) with minimum state variables, rational expectations solution (\ref{sol}) is expectationally stable if $\Lambda<1$. It is not expectationally stable  if $\Lambda>1$.  
\end{proposition}



\clearpage

\section{Data and Model Description}
\label{app:data}

\begin{table}[ht]
\centering
\begin{tabular}{ll p{6.5cm}}
\toprule 
Observable & FRED Series & Construction \\
\midrule
$y^\text{obs}$      & GDPC1, OPHNFB              &
    $\ln(\text{Real GDP}/\text{Labor productivity})$; demeaned \\[3pt]
$d^\text{obs}$      & DPSACBW027SBOG, OPHNFB     &
    $\ln(\text{Total deposits}/\text{Labor productivity})$; demeaned \\[3pt]
$m^\text{obs}$      & WCURCIR, OPHNFB             &
    $\ln(\text{Currency in circulation}/\text{Labor productivity})$; demeaned \\[3pt]
$\pi^\text{obs}$    & GDPDEF                      &
    $\ln(\text{GDPDEF}_t/\text{GDPDEF}_{t-1})$; demeaned \\[3pt]
$r_\ell^\text{obs}$ & MPRIME                      &
    $\ln[(1 + \text{Prime rate}/100)^{0.25}]$; demeaned \\[3pt]
$r_n^\text{obs}$    & FEDFUNDS                    &
    $\ln[(1 + \text{Fed funds rate}/100)^{0.25}]$; demeaned \\[3pt]
$\tau^\text{obs}$   & FGRECPT, GDP                &
    $\ln(\text{Federal tax receipts}/\text{Nominal GDP})$; demeaned \\
\bottomrule
\end{tabular}
\caption{All series from the FRED database (Federal Reserve Bank of St.\ Louis).  OPHNFB is the nonfarm business sector output per hour (labor productivity index).  Interest rates are converted from annualized percent to quarterly gross rates before taking logs.  The term ``mean'' denotes the full-sample arithmetic mean of the transformed series over 2009Q1--2025Q4.  Weekly and monthly series are averaged to quarterly frequency.}
\label{tab:data}
\end{table}

\begin{table}[h]
\centering
\begin{tabular}{llcc}
\toprule
Symbol & Description & Value \\
\midrule
$\beta$      & Household discount factor        & 0.9967 \\
$\chi$       & Housing-service curvature        & 1.0    \\
$\psi$       & Inverse Frisch elasticity        & 1.3    \\
$\theta$     & Goods demand elasticity          & 6.0    \\
$\nu$        & Deposit CES elasticity           & 0.25   \\
$x_n$        & Loan portfolio share             & 0.95   \\
$\phi_n$     & Loan management cost             & 0.0005 \\
$\phi_p$     & Rotemberg adjustment cost        & 50     \\
$\omega$     & Bond share                       & 0.0    \\
$b_0$        & Steady-state debt target         & 1.0    \\
$z$          & Long-run labor productivity      & 1.3    \\
$g$      & Steady-state govt.\ spending     & 0.20   \\
$\tau$    & Steady-state tax rate            & $\{0.18,\, 0.20\}$   \\
$\varphi_b$  & Fiscal feedback to debt          & 0.05   \\
$\rho_\tau$  & Tax persistence                  & 0.50   \\
$\rho_a$     & TFP persistence                  & 0.95   \\
$\rho_{xa}$  & Deposit-technology persistence   & 0.95   \\
$\rho_{va}$  & Preference shock persistence     & 0.0    \\
$\rho_g$     & Govt.\ spending persistence      & 0.0    \\
$\pi$     & Inflation target (gross qtrly.)  & $1.02^{1/4}$ \\
$\underline{r}^{\,n}$ & ZLB floor (gross rate) & 1.0    \\
\bottomrule
\end{tabular}
\caption{Calibrated Parameters}
\label{tab:calib}
\end{table}


\clearpage
%===========================================================
\section{Deficit Case: Full Results ($\tau = 0.18$)}
\label{app:tau018}
%===========================================================

\subsection{Robustness across $\varphi_b$}

\begin{table}[h]
\centering
\small
\begin{tabular}{llcccc}
\toprule
 & & \multicolumn{3}{c}{Posterior mean} \\
\cmidrule(lr){3-5}
Symbol & Description & $\varphi_b\!=\!0.05$ & $\varphi_b\!=\!0.10$ & $\varphi_b\!=\!0.15$ \\
\midrule
$v^a$    & Preference scale               & 2.533 & 2.428 & 2.555 \\
$\eta$   & Labor preference               & 1.456 & 1.427 & 1.434 \\
$b_\pi$  & Inflation response (Taylor)    & 2.113 & 2.122 & 2.121 \\
$b_y$    & Output response (Taylor)       & 0.273 & 0.277 & 0.280 \\
$a_i$    & Policy-rate level weight       & 0.488 & 0.500 & 0.501 \\
$b_r$    & MP shock persistence           & 0.360 & 0.357 & 0.357 \\
$\rho_z$ & Labor-productivity persistence & 0.468 & 0.443 & 0.454 \\
\bottomrule
\multicolumn{4}{l}{\parbox{0.82\linewidth}{\footnotesize Log MDD (mod. harmonic mean): 499.7 / 500.1 / 500.4 for $\varphi_b = 0.05 / 0.10 / 0.15$.}}
\end{tabular}
\caption{Posterior Means across $\varphi_b$ ($\tau = 0.18$)}
\label{tab:posteriors_all_018}
\end{table}

\clearpage
\subsection{Steady-State Values ($\tau = 0.18$, $\varphi_b = 0.05$)}

\begin{table}[h]
\centering
\small
\begin{tabular}{llcc}
\toprule
Symbol & Description & Value & Ann.\ (\%) \\
\midrule
\multicolumn{4}{l}{\textit{Prices and nominal rates}} \\[2pt]
$\pi$          & Gross quarterly inflation    & 1.01211 & 4.93 \\
$r^\ell$       & Gross loan rate              & 1.01546 & 6.33 \\
$r^d$          & Gross deposit rate           & 1.00804 & 3.26 \\
$r^n$          & Gross policy rate            & 1.00751 & 3.04 \\
$w$            & Real wage                    & 1.08333 & ---  \\[4pt]
\multicolumn{4}{l}{\textit{Real quantities}} \\[2pt]
$y$    & Output                         & 1.06099  & --- \\
$c$    & Consumption                    & 0.86099  & --- \\
$h$    & Total labor                    & 0.89466  & --- \\
$d$    & Real deposits                  & 19.185   & --- \\
$n$    & Bank reserves                  &  5.062   & --- \\
$m$    & Monetary aggregate             &  5.120   & --- \\
$\tau$ & Taxes                          &  0.180   & --- \\
\bottomrule
\multicolumn{4}{l}{\parbox{0.82\linewidth}{\footnotesize The fiscal deficit ($g-\tau = 0.02$) generates steady-state inflation of 4.93\% annualised, well above the 2\% target.}}
\end{tabular}
\caption{Steady State at Posterior Means ($\tau = 0.18$, $\varphi_b = 0.05$)}
\label{tab:ss_018}
\end{table}

%===========================================================
\section{Deficit Case: Full Results ($\tau = 0.20$)}
\label{app:tau020}
%===========================================================
\begin{table}[h]
\centering
\small
\begin{tabular}{llcc}
\toprule
Symbol & Description & Value & Ann.\ (\%) \\
\midrule
\multicolumn{4}{l}{\textit{Prices and nominal rates}} \\[2pt]
$\pi$          & Gross quarterly inflation    & 1.00445 & 1.79 \\
$r^\ell$       & Gross loan rate              & 1.00777 & 3.15 \\
$r^d$          & Gross deposit rate           & 1.00149 & 0.60 \\
$r^n$          & Gross policy rate            & 1.00391 & 1.57 \\
$w$            & Real wage                    & 1.08333 & ---  \\[4pt]
\multicolumn{4}{l}{\textit{Real quantities}} \\[2pt]
$y$    & Output                         & 1.08766  & --- \\
$c$    & Consumption                    & 0.88766  & --- \\
$h$    & Total labor                    & 0.91806  & --- \\
$h^g$  & Labor in goods production      & 0.83666  & --- \\
$h^d$  & Labor in deposit management    & 0.07834  & --- \\
$h^n$  & Labor in loan management       & 0.00307  & --- \\
$\mu$  & Marginal utility composite     & 0.96938  & --- \\[4pt]
\multicolumn{4}{l}{\textit{Financial and fiscal variables}} \\[2pt]
$d$    & Real deposits                  & 20.619  & --- \\
$n$    & Bank reserves                  &  6.132  & --- \\
$m$    & Monetary aggregate             &  6.155  & --- \\
$b$    & Government bonds               &  1.000  & --- \\
$g$    & Government spending            &  0.200  & --- \\
$\tau$ & Taxes                          &  0.200  & --- \\
\bottomrule
\multicolumn{4}{l}{\parbox{0.82\linewidth}}
\end{tabular}
\caption{Steady State at Posterior Means ($\tau = 0.20$, $\varphi_b = 0.05$)}
\label{tab:ss_020}
\end{table}

\begin{table}[h]
\centering
\small
\begin{tabular}{llcccc}
\toprule
 & & \multicolumn{3}{c}{Posterior mean} \\
\cmidrule(lr){3-5}
Symbol & Description & $\varphi_b\!=\!0.05$ & $\varphi_b\!=\!0.10$ & $\varphi_b\!=\!0.15$ \\
\midrule

$b_\pi$  & Inflation response (Taylor)    & 2.009 & 1.936 & 2.005 \\
$b_y$    & Output response (Taylor)       & 0.348 & 0.289 & 0.399 \\
$a_i$    & Policy-rate level weight       & 0.457 & 0.682 & 0.534 \\
$b_r$    & MP shock persistence           & 0.404 & 0.383 & 0.394 \\
$v_a$    & Preference scale               & 2.239 & 2.040 & 2.093 \\
$\eta$   & Labor preference               & 1.430 & 1.344 & 1.313 \\
$\rho_z$ & Labor-productivity persistence & 0.463 & 0.529 & 0.464 \\
\bottomrule
\multicolumn{4}{l}{\parbox{0.82\linewidth}}
\end{tabular}
\caption{Posterior means for $\varphi_b = 0.05 / 0.10 / 0.15$. ($\tau = 0.20$)}
\label{tab:posteriors_all_020}
\end{table}

\end{document}

